% GENERATED FILE. Edit canonical lessons, then run npm run generate:books.
% canonical-source-hash: fnv1a64:7a7f619097291716
% canonical-lessons: SA-C154-yadi, SA-C154-cet, SA-C154-tathapi, SA-C154-yadyapi, SA-C154-athava

\chapter{If, Provided that, Even so, Although, Or else}
\label{ch:sa-if-provided-that-even-so-although-or-else}
\hlchaptermodality{154}

\begin{chapteropening}
\textbf{By the end of this chapter:} \emph{I can say if, provided that, even so, although and or else.}

Complete the last lesson of chapter 154: I can say if, provided that, even so, although and or else.
\end{chapteropening}

\section[\texorpdfstring{\sk{यदि} (yadi)}{यदि (yadi)}]{\sk{यदि} (yadi) — if}

\label{lesson:SA-C154-yadi}



\begin{quote}
Before the new one: say the Sanskrit for all around, then the Sanskrit for on all sides.
\end{quote}

\subsection*{You'll want to know: \sk{यदि}}
\textbf{\sk{यदि}} — \emph{yadi} — \textquotedblleft{}if\textquotedblright{}.

\begin{grammarlens}[title={What you've built}]
One more word for this chapter.
\end{grammarlens}

\subsection*{Your turn}
\begin{itemize}
\raggedright
  \item \emph{Say it:} \emph{yadi}
  \item \emph{Say it:} \emph{yadi}, once more
  \item \emph{Say it:} say \emph{yadi}
  \item \emph{Recall:} say the Sanskrit for all around, then the Sanskrit for on all sides, then say \emph{yadi} again
\end{itemize}

\begin{tcolorbox}[breakable,colback=teal!4,colframe=teal!35!black,title={Before you move on}]
What does \sk{यदि} mean? (\textquotedblleft{}If\textquotedblright{}.) Say it once more.
\end{tcolorbox}
\section[\texorpdfstring{\sk{चेत्} (cet)}{चेत् (cet)}]{\sk{चेत्} (cet) — provided that, if}

\label{lesson:SA-C154-cet}



\begin{quote}
Before the new one: say the Sanskrit for on all sides, then the Sanskrit for if.
\end{quote}

\subsection*{You'll want to know: \sk{चेत्}}
\textbf{\sk{चेत्}} — \emph{cet} — \textquotedblleft{}provided that, if\textquotedblright{}.

\begin{grammarlens}[title={What you've built}]
One more word for this chapter.
\end{grammarlens}

\subsection*{Your turn}
\begin{itemize}
\raggedright
  \item \emph{Say it:} \emph{cet}
  \item \emph{Say it:} \emph{cet}, once more
  \item \emph{Say it:} say \emph{cet}
  \item \emph{Recall:} say the Sanskrit for on all sides, then the Sanskrit for if, then say \emph{cet} again
\end{itemize}

\begin{tcolorbox}[breakable,colback=teal!4,colframe=teal!35!black,title={Before you move on}]
What does \sk{चेत्} mean? (\textquotedblleft{}Provided that, if\textquotedblright{}.) Say it once more.
\end{tcolorbox}
\section[\texorpdfstring{\sk{तथापि} (tathāpi)}{तथापि (tathāpi)}]{\sk{तथापि} (tathāpi) — even so, still}

\label{lesson:SA-C154-tathapi}



\begin{quote}
Before the new one: say the Sanskrit for if, then the Sanskrit for provided that.
\end{quote}

\subsection*{You'll want to know: \sk{तथापि}}
\textbf{\sk{तथापि}} — \emph{tathāpi} — \textquotedblleft{}even so, still\textquotedblright{}.

\begin{grammarlens}[title={What you've built}]
One more word for this chapter.
\end{grammarlens}

\subsection*{Your turn}
\begin{itemize}
\raggedright
  \item \emph{Say it:} \emph{tathāpi}
  \item \emph{Say it:} \emph{tathāpi}, once more
  \item \emph{Say it:} say \emph{tathāpi}
  \item \emph{Recall:} say the Sanskrit for if, then the Sanskrit for provided that, then say \emph{tathāpi} again
\end{itemize}

\begin{tcolorbox}[breakable,colback=teal!4,colframe=teal!35!black,title={Before you move on}]
What does \sk{तथापि} mean? (\textquotedblleft{}Even so, still\textquotedblright{}.) Say it once more.
\end{tcolorbox}
\section[\texorpdfstring{\sk{यद्यपि} (yadyapi)}{यद्यपि (yadyapi)}]{\sk{यद्यपि} (yadyapi) — although}

\label{lesson:SA-C154-yadyapi}



\begin{quote}
Before the new one: say the Sanskrit for provided that, then the Sanskrit for even so.
\end{quote}

\subsection*{You'll want to know: \sk{यद्यपि}}
\textbf{\sk{यद्यपि}} — \emph{yadyapi} — \textquotedblleft{}although\textquotedblright{}.

\begin{grammarlens}[title={What you've built}]
One more word for this chapter.
\end{grammarlens}

\subsection*{Your turn}
\begin{itemize}
\raggedright
  \item \emph{Say it:} \emph{yadyapi}
  \item \emph{Say it:} \emph{yadyapi}, once more
  \item \emph{Say it:} say \emph{yadyapi}
  \item \emph{Recall:} say the Sanskrit for provided that, then the Sanskrit for even so, then say \emph{yadyapi} again
\end{itemize}

\begin{tcolorbox}[breakable,colback=teal!4,colframe=teal!35!black,title={Before you move on}]
What does \sk{यद्यपि} mean? (\textquotedblleft{}Although\textquotedblright{}.) Say it once more.
\end{tcolorbox}
\section[\texorpdfstring{\sk{अथवा} (athavā)}{अथवा (athavā)}]{\sk{अथवा} (athavā) — or else, or}

\label{lesson:SA-C154-athava}



\begin{quote}
Before the new one: say the Sanskrit for even so, then the Sanskrit for although.
\end{quote}

\subsection*{You'll want to know: \sk{अथवा}}
\textbf{\sk{अथवा}} — \emph{athavā} — \textquotedblleft{}or else, or\textquotedblright{}.

\begin{grammarlens}[title={What you've built}]
One more word for this chapter.
\end{grammarlens}

\subsection*{Your turn}
\begin{itemize}
\raggedright
  \item \emph{Say it:} \emph{athavā}
  \item \emph{Say it:} \emph{athavā}, once more
  \item \emph{Say it:} say \emph{athavā}
  \item \emph{Recall:} say the Sanskrit for even so, then the Sanskrit for although, then say \emph{athavā} again
\end{itemize}

\begin{tcolorbox}[breakable,colback=teal!4,colframe=teal!35!black,title={Before you move on}]
What does \sk{अथवा} mean? (\textquotedblleft{}Or else, or\textquotedblright{}.) Say it once more.
\end{tcolorbox}
